\documentclass[12pt]{article}
\usepackage[utf8]{inputenc}
\usepackage[T1]{fontenc}
\usepackage{hyperref}
\title{Sanie Language Resources for K-12 Curriculum Development}
\author{Compiled by Claude Sonnet 5.5 (claude-sonnet-5-5) for Andr\'as Kornai}
\date{}
\begin{document}
\maketitle

\section{Introduction}
Sanie (ISO 639-3 ysy; Glottocode sani1265) is a Loloish language of central Yunnan, China. Wikipedia reports about 17,320 ethnic Sanie in 1998, of whom only about 8,000 spoke the language fluently. Some sources class it as a dialect of Sani (ysn), a different language that must not be confused with Sanie. Endangered Languages lists Sanie as severely endangered. Very little material exists for it, and nothing found suggests it is taught in schools.

\section{Primary Educational Resources}
No primers, textbooks or school programmes for Sanie turned up in the searches. A Sanie pinyin orthography is reported by Wikipedia, citing Bradley (2005), to have been developed for language maintenance. The searches did not locate its materials.

\section{Linguistic and Reference Materials}
\begin{itemize}
\item Wikipedia, Sanie language, with dialect information (Bradley 2005): \url{https://en.wikipedia.org/wiki/Sanie_language}
\item Endangered Languages Project entry for ysy: \url{https://www.endangeredlanguages.com/lang/ysy}
\item Wikipedia, Sani language (the related but distinct ysn): \url{https://en.wikipedia.org/wiki/Sani_language}
\item Omniglot, Sani (Nip-dop), the related Sani language, which has Yi-script and Latin writing: \url{https://www.omniglot.com/writing/sani.htm}
\end{itemize}
The main academic sources named in the search results are David Bradley's work (2005) and Ma Xueliang's 1951 study of Sani. Their full texts were not found.

\section{Religious and Community Resources}
\begin{itemize}
\item Asia Harvest people-group profile for the Sanie: \url{https://people-groups.asiaharvest.org/people-group-profiles/yunnan/sanie.pdf}
\end{itemize}
No community-run language organization was found.

\section{Digital and Audio Resources}
\begin{itemize}
\item Global Recordings Network page for ysy, which lists no recordings at the time of the search: \url{https://globalrecordings.net/en/language/ysy}
\end{itemize}

\section{Curriculum Development Context}
Few resources exist, and children in nearby villages are reported to be shifting to other languages. Any curriculum would therefore start from documentation work with older speakers, using the Sanie pinyin orthography or the Yi script, and would draw on the Sani materials only with caution. The language is not known to be taught in schools.
\end{document}